
%custom commands for folder structure - please use :)
\newcommand{\meta}{meta}
\newcommand{\ana}{chapters/analysis}
\newcommand{\conc}{chapters/conclusion}
\newcommand{\intro}{chapters/introduction}
\newcommand{\fund}{chapters/fundamentals}
\newcommand{\conf}{config}
\newcommand{\appx}{appendix}
\newcommand{\bibl}{bibliography}
\newcommand{\imagedir}{img}

\documentclass[fontsize=12pt,BCOR=5mm,DIV=12,
               parskip=half,
               listof=entryprefix,paper=a4,toc=bibliography,toc=listof,pointlessnumbers
               %,headinclude=on,footinclude=off,
               %,plainheadtopline=false,plainheadsepline=false,plainfootsepline=false,plainfootbotline=false
               ]{scrreprt}
               
\makeindex

% (!!) Elementare Pakete, Konfigurationen und Definitionen werden geladen
% !TEX root =  master.tex
%      HYPERREF

%%%%%%%%%%%%%%%%%%%%%%%%%%%%%%%%%%%%%%%%%%%%%%%%%%%%%%%%%%
%	ANLEITUNG: 
%
% Passen Sie alle Stellen im Dokument an, die mit 
% @stud markiert sind
%
%%%%%%%%%%%%%%%%%%%%%%%%%%%%%%%%%%%%%%%%%%%%%%%%%%%%%%%%%%

\usepackage{makeidx}         % allows index generation
\usepackage{listings}	%Format Listings properly
\usepackage{lipsum}    %Blindtext
\usepackage{graphicx} % use various graphics formats
\usepackage[german]{varioref} 	% nicer references \vref
\usepackage{caption}	%better Captions
\usepackage{booktabs} %nicer Tabs
\usepackage{array}
\usepackage{chngcntr}
\usepackage[hidelinks=true]{hyperref} % keine roten Markierungen bei Links
\usepackage{fnpct} % Correct superscripts 
\usepackage[T1]{fontenc}
\usepackage[utf8]{inputenc}
\usepackage{calc} % Used for extra space below footsepline
\usepackage{acronym}
\usepackage{algorithm}
\usepackage{algpseudocode}
\usepackage{setspace}

%
% @stud
%
%	FONT SELECTION: Entweder 1) Latin Modern oder 2) Times / Helvetica
\usepackage{lmodern}             % 1) Latin modern font
%\usepackage{mathptmx}           % 2) Helvetica / Times New Roman fonts (2 lines)
%\usepackage[scaled=.92]{helvet} % 2) Helvetica / Times New Roman fonts (2 lines)

%
% @stud
%
%	LANGUAGE SETTINGS
\usepackage[ngerman]{babel} 	        % german language
\usepackage[german=quotes]{csquotes} 	% correct quoting using \enquote{}
%\usepackage[english]{babel}          % english language
%\usepackage{csquotes} 	              % correct quoting using \enquote{}

%
% @stud
%
% Uncomment the following lines to support hard URL breaks in bibliography 
%\apptocmd{\UrlBreaks}{\do\f\do\m}{}{}
%\setcounter{biburllcpenalty}{9000}% Kleinbuchstaben
%\setcounter{biburlucpenalty}{9000}% Großbuchstaben

%
% @stud
%
%	FOOTNOTES: Count footnotes over chapters
%1 \counterwithout{footnote}{chapter}

%	ACRONYMS
\makeatletter
\@ifpackagelater{acronym}{2015/03/20}
{\renewcommand*{\aclabelfont}[1]{\textbf{{\acsfont{#1}}}}}{}
\makeatother

%	LISTINGS
\renewcommand{\lstlistingname}{Quelltext}
\renewcommand{\lstlistlistingname}{Quelltextverzeichnis}
\lstset{numbers=left,
	numberstyle=\tiny,
	captionpos=b,
	basicstyle=\ttfamily\small}

%	ALGORITHMS
\renewcommand{\listalgorithmname}{Algorithmenverzeichnis }
\floatname{algorithm}{Algorithmus}

%		PAGE HEADER / FOOTER
%	    Warning: There are some redefinitions throughout the master.tex-file!  DON'T CHANGE THESE REDEFINITIONS!
\RequirePackage[automark]{scrlayer-scrpage}
%alternatively with separation lines: \RequirePackage[automark,headsepline,footsepline]{scrlayer-scrpage}

%\renewcommand*{\pnumfont}{\upshape\sffamily}
%\renewcommand*{\headfont}{\upshape\sffamily}
%\renewcommand*{\footfont}{\upshape\sffamily}

\renewcommand{\chaptermarkformat}{}
\RedeclareSectionCommand[beforeskip=0pt]{chapter}
\clearscrheadfoot

%\ifoot[\rule{0pt}{\ht\strutbox+\dp\strutbox}DHBW Mannheim]{\rule{0pt}{\ht\strutbox+\dp\strutbox}DHBW Mannheim}
\ofoot[\rule{0pt}{\ht\strutbox+\dp\strutbox}\pagemark]{\rule{0pt}{\ht\strutbox+\dp\strutbox}\pagemark}
\ohead{\headmark}

\newcommand{\TitelDerArbeit}[1]{\def\DerTitelDerArbeit{#1}\hypersetup{pdftitle={#1}}}
\newcommand{\AutorDerArbeit}[1]{\def\DerAutorDerArbeit{#1}\hypersetup{pdfauthor={#1}}}
\newcommand{\Firma}[1]{\def\DerNameDerFirma{#1}}
\newcommand{\Kurs}[1]{\def\DieKursbezeichnung{#1}}
\newcommand{\Studiengangsleiter}[1]{\def\DerStudiengangsleiter{#1}}
\newcommand{\Bearbeitungszeitraum}[1]{\def\DerBearbeitungszeitraum{#1}}
\newcommand{\Abgabedatum}[1]{\def\DasAbgabedatum{#1}}
\newcommand{\Matrikelnummer}[1]{\def\DieMatrikelnummer{#1}}
\newcommand{\Studienrichtung}[1]{\def\DieStudienrichtung{#1}}
\newcommand{\ArtDerArbeit}[1]{\def\DieArtDerArbeit{#1}}
\newcommand{\Literaturverzeichnis}{Literaturverzeichnis}

\newcommand{\settingBibFootnoteCite}{
	\setlength{\bibparsep}{\parskip}		  % Add some space between biblatex entries in the bibliography
	\addbibresource{\bibl/bibliography.bib}	    % Add file bibliography.bib as biblatex resource
	\DefineBibliographyStrings{ngerman}{andothers = {{et\,al\adddot}},}
	\AdaptNoteOpt\footcite\multfootcite   % Will add  separators if footcite is called multiple consecutive times 
	\AdaptNoteOpt\autocite\multautocite   % Will add  separators if autocite is called multiple consecutive times
}

\newcommand{\setTitlepage}{
	% !TEX root =  master.tex
\begin{titlepage}
	\begin{minipage}{\textwidth}
		\vspace{-2cm}
 \includegraphics{\imagedir/logo.jpg}
	\end{minipage}
	\vspace{1em}
	%\sffamily
	\begin{center}
		{\textsf{\large Duale Hochschule Baden-W\"urttemberg Mannheim}}\\[4em]
		{\textsf{\textbf{\large{\DieArtDerArbeit}arbeit}}}\\[6mm]
		{\textsf{\textbf{\Large{}\DerTitelDerArbeit}}} \\[1.5cm]
		{\textsf{\textbf{\large{}Studiengang Informatik}}\\[6mm]
		\textsf{\textbf{Studienrichtung \DieStudienrichtung}}}\vspace{10em}

		\begin{minipage}{\textwidth}
			\begin{tabbing}
				Wissenschaftliche: \hspace{0.85cm}\=\kill
				Verfasser: \> Max Stubenbord \\ \>  Martin Pretz \\ \> Abdalroman Alsalkhadi\\[1.5mm]
				Matrikelnummer: \> \DieMatrikelnummer \\[1.5mm]
				Kurs: \> \DieKursbezeichnung \\[1.5mm]
				Studiengangsleiter: \> \DerStudiengangsleiter \\[1.5mm]
				Bearbeitungszeitraum: \> \DerBearbeitungszeitraum\\[1.5mm]
				%		alternativ:\\[1.5mm]
				%		Eingereicht: \> \DasAbgabedatum	
			\end{tabbing}
		\end{minipage}
	\end{center}
\end{titlepage}
	\pagenumbering{roman} % Römische Seitennummerierung
	\normalfont
}

%
% @stud
%
\newcommand{\settingLists}{
	%	Inhaltsverzeichnis
	\tableofcontents
	%	Abbildungsverzeichnis
	\listoffigures
	%	Tabellenverzeichnis
	\listoftables
	%	Listingsverzeichnis / Quelltextverzeichnis
	\lstlistoflistings
	% Algorithmenverzeichnis
	\listofalgorithms
}

\newcommand{\initializeText}{
	\clearpage
	\ihead{\chaptername~\thechapter} % Neue Header-Definition
	\pagenumbering{arabic}           % Arabische Seitenzahlen
}

\newcommand{\initializeBibliography}{
	\ihead{}
	\printbibliography[title=\Literaturverzeichnis]
	\cleardoublepage
}

\newcommand{\initializeAppendix}{
	\appendix
	\ihead{\appendixname~\thechapter}
}



%%%%%%%%%%%%%%%%%%%%%%%%%%%%
%
% @stud
%
%	SCHRIFTART (Schrift mit oder ohne Serifen im gesamten Text) 
%
% mit Serifen
%\addtokomafont{disposition}{\rmfamily}
%\renewcommand*{\familydefault}{\rmdefault}
%
% ohne Serifen (default)
%\addtokomafont{disposition}{\sffamily}
%
%%%%%%%%%%%%%%%%%%%%%%%%%%%%

%%%%%%%%%%%%%%%%%%%%%%%%%%%%
%
% @stud
%
% PERSÖNLICHE ANGABEN (BITTE VOLLSTÄNDIG EINGEBEN zwischen den Klammern: {...})
%
\ArtDerArbeit{Projekt} % "Bachelor" oder "Projekt" wählen
\TitelDerArbeit{Iwas mit Rust}
\AutorDerArbeit{Max Stubenbord,  Martin Pretz, Abdalrohman Alsalkhadi}
\Kurs{TINF18AI1}
\Studienrichtung{Angewandte Informatik}
\Matrikelnummer{5379506, 7060026, 1019128}
\Studiengangsleiter{Alexander Kühn}
\Bearbeitungszeitraum{15.11.2020 -- 20.11.2020}
\Abgabedatum{20.11.2020}
%
%%%%%%%%%%%%%%%%%%%%%%%%%%%%

%%%%%%%%%%%%%%%%%%%%%%%%%%%%
%
% @stud
%
%	BIBLIOGRAPHY (@stud: Bibliographie-Stil wählen - Position und Indizierung)
%
% Auswahl zwischen: IEEE Style, ALPHABETIC Style, HARVARD Style, AUTHOR-YEAR Style 
%
% (oder eigenen zulässigen Stil wählen) 
%

% Position des Zitats
%




\newcommand{\position}{inline} 
%\newcommand{\position}{footnote}

% Indizierung des Zitats
%
% 1) NUMERIC Style - e. g. [12]
\newcommand{\indextype}{numeric} 
%
% 2) ALPHABETIC Style - e. g. [AB12]
%\newcommand{\indextype}{alphabetic} 
%
% 3) IEEE Style - numeric kind of style 
%\newcommand{\indextype}{ieee} 
%
% 4) HARVARD Style 
%\newcommand{\indextype}{apa} 
%
% 5) CHICAGO Style 
%\newcommand{\indextype}{authoryear}
%
%%%%%%%%%%%%%%%%%%%%%%%%%%%%

\renewcommand*{\familydefault}{\sfdefault}

\usepackage[backend=biber, autocite=\position, style=\indextype]{biblatex} 	

\settingBibFootnoteCite

\newcommand{\abs}{\par\vskip 0.2cm\goodbreak\noindent}
\newcommand{\nl}{\par\noindent}
\newcommand{\mcl}[1]{\mathcal{#1}}
\newcommand{\nowrite}[1]{}
\newcommand{\NN}{{\mathbb N}}


\makeindex

\begin{document}

\setTitlepage


% !TEX root =  master.tex
\clearpage
\chapter*{Ehrenwörtliche Erklärung}

% Wird die folgende Zeile auskommentiert, erscheint die ehrenwörtliche
% Erklärung im Inhaltsverzeichnis.

% \addcontentsline{toc}{chapter}{Ehrenwörtliche Erklärung}
Wir versichern hiermit, dass wir die vorliegende Arbeit
mit dem Thema: \textit{\DerTitelDerArbeit} selbstständig verfasst und keine anderen als die angegebenen Quellen und
Hilfsmittel benutzt habe. Ich versichere zudem,
dass die eingereichte elektronische Fassung mit der gedruckten Fassung übereinstimmt.

\vspace{3cm}
Ort, Datum \hfill \DerAutorDerArbeit


% !TEX root =  master.tex
\chapter*{Sperrvermerk}

\begin{center}
\fbox{
		\begin{minipage}{33em}
			\textbf{Ein Sperrvermerk sollte nur bei berechtigtem Bedarf gesetzt werden!\\[10pt] 
				Beachten Sie, dass mit Sperrvermerk	versehene Arbeiten nicht für weitere wissenschaftliche Zwecke 
				außerhalb des Firmenkontextes oder zur Publikation verwendet werden dürfen.\\[10pt]
				Wir empfehlen, wenn m\"oglich, auf den Sperrvermerk zu verzichten.\\[10pt]
				Besprechen Sie diese Problematik mit Ihrer Firma!}
		\end{minipage}
}
\end{center}

(Mustertext) Der Inhalt dieser Arbeit darf weder als Ganzes noch in Auszügen Personen außerhalb des Prüfungsprozesses 
und des Evaluationsverfahrens zugänglich gemacht werden, sofern keine anders lautende Genehmigung der Ausbildungsstätte vorliegt. 

\cleardoublepage


% !TEX root =  master.tex
\chapter*{Danksagung}

Hier können Sie eine Danksagung schreiben. 




% !TEX root =  master.tex
\chapter*{Kurzfassung (Abstract)}

Hier können Sie die Kurzfassung (engl.~Abstract) der Arbeit schreiben. Beachten Sie dabei die Hinweise zum Verfassen der Kurzfassung.


\settingLists

% !TEX root =  master.tex
\clearpage
\chapter*{Abkürzungsverzeichnis}	
\addcontentsline{toc}{chapter}{Abkürzungsverzeichnis}

\begin{acronym}[XXXXXXX]
	\acro{ad}[AD]{Archiv für Diplomatik, Schriftgeschichte, Siegel- und Wappenkunde}
	\acro{BMBF}{Bundesministerium für Bildung und Forschung}	
	\acro{DHBW}{Duale Hochschule Baden-Württemberg}
	\acro{ecu}[ECU]{European Currency Unit}
	\acro{eu}[EU]{Europäische Union}
	\acro{RDBMS}[RDBMS]{Relational Database Management System}
\end{acronym}


\initializeText
\onehalfspacing



\chapter{Einleitung}

\nocite{*}
Rust ist eine moderne Multiparadigmen-Systemprogrammiersprache, die sich auf Sicherheit, Geschwindigkeit und Parallelität konzentriert. Sie entstand 2006 durch ein persönliches Projekt des Mozilla-Mitarbeiters Graydon Hoare und entwickelte sich im Laufe von fünf Jahren zu einem Open-Source-Projekt, das von Mozilla Research entwickelt wird und von einer großen Community von freiwilligen Mitwirkenden unterstützt wird. Traurigerweise ist die übliche Auffassung, dass man nicht alle gewünschten Eigenschaften in einer Programmiersprache haben kann. Sprachen wie Java bieten eine starke Typ- und Speichersicherheit und verhindern damit eine Vielzahl von Speicherfehlern, allerdings geht das auf Kosten der Ressourcenkontrolle. Infolgedessen besteht für viele Anwendungen zur Systemprogrammierung die einzige realistische Option darin, eine Sprache wie \texttt{C} oder \texttt{C++} zu verwenden, die zwar eine fein granulierte Kontrolle über die Speicherverwaltung bietet, jedoch ist diese Kontrolle mit hohen Kosten verbunden. Beispielsweise berichtete Microsoft kürzlich, dass 70\% der Sicherheitslücken, die sie beheben, auf Verletzungen der Speichersicherheit zurückzuführen sind, also genau die Art von Verletzungen, die bei starken Typsystemen ausgeschlossen werden sollten. Ebenso berichtet Mozilla, dass die überwiegende Mehrheit der kritischen Fehler, die sie in Firefox finden, speicherbezogen sind. Wenn es nur einen Weg gäbe, irgendwie das Beste aus beiden Welten zu vereinen, sodass eine sichere Systemprogrammiersprache entsteht, dann würde man versuchen, so eine Sprache wie Rust zu entwickeln. Rust enstand aus der Note heraus. Sie ist die erste von der Industrie unterstützte Programmiersprache, die den seit langem bestehenden Kompromiss im Entwurf von Programmiersprachen zwischen den Sicherheitsgarantien höherer Programmiersprachen (wie Java) und der Kontrolle über die Speicherverwaltung durch low-level-Programmiersprachen (wie \texttt{C} und \texttt{C++}) überwindet. Rust geht diese Herausforderung mit Hilfe eines starken, neuartigen Typsystems an, das auf den Ideen der Eigentümerschaft (Ownership) and des Ausleihen (Borrowing) basiert und die statische Mutation eines gemeinsamen Zustandes verbietet. Dieser Ansatz ermöglicht es, die üblichen Fallstricke der Systemprogrammierung wie Data Races zur Kompilierzeit zu erkennen und zu verhindern. 

Das Ownership-System legt für jeden Wert eine klare Lebensdauer fest, wodurch die Verwendung eines Garbage-Collectors überflüssig wird und solide, aber flexible Schnittstellen für die Verwaltung anderer Arten von Ressourcen wie Sockets und Datei-Handles ermöglicht werden. Durch \glqq Moves\grqq{} können Werte von einem Eigentümer auf einen anderen übertragen werden, und durch Ausleihen kann ein Code einen Wert vorübergehend verwenden, ohne dass sich dies auf das Eigentum auswirkt. 
Die Kontrolle über die Speicherverwaltung ist ein sehr wichtiger Aspekt der systemnahen Programmierung, besonders für Anwendungen wie Web-Browser, Betriebssysteme oder Game-Engines, bei denen Leistungs- oder Ressourcenbeschränkungen ein Hauptanliegen sind. Diese Art der Kontrolle wird von Rust unterstützt, womit die Möglichkeit gegeben wird, die Darstellung von Daten auf Byte-Ebene zu bestimmen sowie einen Zugang zum \glqq Bare Metal\grqq{} zu schaffen und die zeitliche und räumliche Nutzung von Programmen durch Low-Level-Programmiertechniken zu optimieren.
Es gibt eine Reihe von Datentypen, deren Implementierungen grundlegend vom gemeinsam veränderlichen Zustand (shared mutable state) abhängen und daher nicht nach der strengen Ownership-Prinzip von Rust typgeprüft werden können. Um solche Datentypen zu unterstützen, setzt Rust auf die vernünftige Verwendung von unsicheren Codes (engl. Unsafe Code), die in sicheren APIs gekapselt sind.
Diese Alleinstellungsmerkmale, über die nur Rust verfügt, sind wohl der Grund, warum Rust von Entwicklern in einer 2020 von der Stack Overflow durchgeführten Umfrage zum fünften Mal in Folge zur beliebtesten Programmiersprache gewählt wurde. 




\chapter{Design}
Zum Design einer Programmiersprache gehören im Allgemeinen viele verschiedene Elemente und Konzepte die dazu beitragen, dass eine Programmiersprache sicher und zuverlässig ist. Aufgrund des erheblichen Umfangs der Elemente und Konzepte und deren Feinheiten, werden in diesem Kapitel lediglich vier ausgewählte Bestandteile des Designs von Rust vorgestellt. \autocite{rust-wiki}
\section{Syntax}
Rust ist syntaktische ähnlich zu der Programmiersprache C. Konkret ermöglicht die Syntax von Rust sowohl die einfache Lesbarkeit des Codes, als auch das ausschließen von grundlegenden Formatierungsfehlern. [1]
Im Endeffekt bedeutet dies, dass der Code aus Blöcken besteht, welche mit geschweiften Klammern begonnen und beendet werden. Dies gilt jedoch nicht für die Bedingungen von Schleifen wie zum Beispiel while oder for sehr wohl aber für deren Schleifeninhalt. Die einzelnen Anweisungen werden dabei stets mit einem Semikolon voneinander getrennt.  [1] [2] [3]
Viele der syntaktischen Merkmale von Rust sind von Grund auf sicher und werden daher intern als "safe" gelistet. Allerdings bietet Rust auch die Möglichkeit für den Entwickler sich von der Sicherheit wegzubewegen und bewusst Code zu schreiben, welcher nach Rust-Kriterien als unsicher klassifiziert wird. Dafür muss allerdings der entsprechende Codeabschnitt im Code explizit in einem Rahmen mit dem Schlüsselwort "unsafe" angegeben werden. [1] [2] [3]
\newpage
\begin{minted}[linenos, fontsize=\small]{rust}
    fn main() {
        struct Viereck {
            x: f32,
            y: f32
        }
        impl Viereck {
            fn flaeche(&self) -> f32 {
                self.x * self.y
            }
        }
    }
\end{minted}
\begin{lstlisting}[caption={Einfaches Struct und zugeh"orige Methode \\Quelle: \autocite{rust-wiki}}, label={lst:struct}]
\end{lstlisting}

\section{Typen}

\section{Memory}

\section{Ownership}
\chapter{Anwendungsbereiche} % (fold)
\label{cha:Anwendungsbereiche}
\section{statisch und dynamisch Typisierte Sprachen}
Die Programmiersprache Rust ist seit vier Jahren in Folge die beliebteste Sprache von Stack Overflow, was darauf hindeutet, dass viele, die die Gelegenheit hatten Rust zu benutzen, sich in diese Sprache verliebt haben. \autocite{so-rust-loved}\\
Ein Indikator, f"ur diese Beliebtheit ist wohl, dass Rust an jenen Stellen ansetzt, welche in anderen Sprachen Entwicklern Bauchschmerzen bringen.\\
In den letzten Jahren wurde die Auseinandersetzung zwischen Programmierern, welche dynamische Sprachsysteme den statischen bevorzugen immer mehr zum Thema. Mit dem Aufstieg von Typescript oder dem feature des typisierten Pythons wurde die statische Typisierung wieder zum Thema in der IT-Branche.\\
Dynamische Typisierung in gr"o"seren Codebasen frustriert Entwickler schon seit l"angerer Zeit. \\
Statisch typisierte Sprachen ermöglichen vom Compiler überprüfte Einschränkungen der Daten und ihres Verhaltens, wodurch kognitiver Overhead und Missverständnisse verringert werden. Somit steigt auch der Anteil von statisch typisierten Projekten wieder. \\
Dies induziert aber keinenfalls, dass alle statischen Systeme gleichwertig zu betrachten sind. Viele statisch typisierte Sprachen kommen mit einem Asterix neben dem Wort \textbf{statisch Typisiert*}.\\
Diese Anmerkung sagt aus, dass diese Sprachen das Konzept von \textbf{NULL} unterst"utzen.
Hierbei kann jeder Wert zwei Zust"ande besitzen: Den vom Entwickler angegebenen Wert oder eben NULL. \\
Wie auch beispielsweise die Programmiersprache Haskell l"ost Rust dieses Problem mit einem optionalem Typ. Dabei wird vom Compiler verlangt, dass der Fall \textbf{None} vom Entwickler behandelt wird.\\
Durch diese Implementierung der Nullverwertung werden systematisch alle gef"urchteten \textbf{Typeerrors} abgefangen und somit auch verhindert.\\
Ein Anwendungsbeispiel in Rust hierf"ur ist eine einfache Begr"u"sung, welche schaut, ob ein Name angegeben wurde.\\
Hier ein kleines Codebeispiel:
\newpage
\begin{minted}[linenos, fontsize=\small]{rust}
    fn greet_user(name: Option<String>) {
    match name {
        Some(name) => println!("Hello there, {}!", name),
        None => println!("Well howdy, stranger!"),
    }
}

 \end{minted}
\begin{lstlisting}[caption={Optionale Typisierung}, label={lst:optionaltypes}]
 \end{lstlisting}

Für viele wird Rust wie schon  in der Einleitung erw"ahnt als Alternative zu den Programmiersprachen wie \texttt{C} oder \texttt{C++} angesehen. Der größte Vorteil, den Rust im Vergleich zu diesen Sprachen bieten kann, ist der \textbf{Borrow-Checker}.\\
Dies ist der Teil des Compilers, welcher dafür verantwortlich ist, dass der Lebenszyklus der Referenzen nicht dem der darauf verweisenden Daten "uberschreitet.\\
Rust ist bestrebt, so viele kostenfreie Abstraktionen wie möglich zu haben - Abstraktionen, die genauso performant sind wie der entsprechend geschriebene Code.\\

Anhand dieses Beispiels wird gezeigt, wie Iteratoren verwendet werden können, um einen Vektor zu erzeugen, der die ersten zehn Quadratzahlen enthält:

\begin{minted}[fontsize=\small]{rust}
    let squares: Vec<_> = (0..10).map(|i| i * i).collect();
\end{minted}
\begin{lstlisting}[caption={Abstrakte Iteration}, label={lst:abstract}]
\end{lstlisting}

Systemprogrammiersprachen haben die implizite Erwartung, dass es sie für immer geben wird. \\
Während einige moderne Entwicklungen nicht diese Langlebigkeit erfordern, wollen viele Unternehmen, dass ihre grundlegende Code-Basis in absehbarer Zukunft immernoch nutzbar sein wird.\\
Rust ist sich dessen bewusst und hat bewusste Design-Entscheidungen im Hinblick auf Rückwärtskompatibilität und Stabilität getroffen; es ist eine Sprache, die für die \enquote{n"achsten 40 Jahre entworfen wurde}.

\section{Ecosystem}
Rust ist gr"o"ser als nur eine Sprachspezifikation und der dazugeh"orige Compiler; viele Aspekte der Erstellung und Wartung von Software werden von Rust priorisiert.\\
Mehrere gleichzeitige Rust Toolchains können über \textbf{Rustup} \autocite{rustup} installiert und verwaltet werden. Rust-Installationen werden mit Cargo geliefert, einem Befehlszeilen-Tool zur Verwaltung von Abhängigkeiten, zur Durchführung von Tests, zur Erstellung von Dokumentation etc. .\\
Da Abhängigkeiten, Tests und Dokumentation standardmäßig verfügbar sind, werden sie in Rust Projekten auch dementsprechend h"aufig eingesetzt. \textbf{Cratesio} \autocite{cratesio} ist die Community-Site zur gemeinsamen Nutzung und Entdeckung von Rust-Bibliotheken. Jede Bibliothek, die auf crates.io veröffentlicht wird, lässt ihre Dokumentation auf \textbf{docsrs} \autocite{docsrs} erstellen und veröffentlichen.\\

Zusätzlich zu den schon vorhandenen Tools hat die Rust-Community viele weitere Tools  entwickelt.\\
Hierzu z"ahlen Benchmarking, Fuzzing und Tests. Diese sind alle leicht zugänglich und werden in Projekten gut genutzt. Zusätzliche Compiler-Lints sind von Clippy erhältlich und automatische idiomatische Formatierung wird von rustfmt bereitgestellt.

Über die technischen Aspekte hinaus hat Rust eine sehr aktive und hilfsbereite Community. Es gibt mehrere offizielle und inoffizielle Möglichkeiten, Hilfe zu erhalten:
\begin{itemize}
    \item  Chat \autocite{rust-discord}
    \item Benutzerforum \autocite{rust-user-forum}
    \item Subreddit \autocite{rust-subreddit}
\end{itemize}

\section{Anwendung}
Rust findet neben Privatem Nutzen auch im kommerziellen Bereich eine Anwendung. Hierbei nutzen Firmen, wie \textbf{npm, tilde, Dropbox} oder \textbf{Yelp} die Sprache in verschiedenen Einsatzbereichen.\\
Npm setzt Rust ein, um \enquote{cpu bound bottlenecks} zu umgehen.
Ein Whitepaper indiziert den Nutzen von Rust bei npm und erl"autert hierbei warum Rust eingesetzt wurde.\\
Im Fazit des Papers wird als einer der Gr"unde, wieso Rust benutzt wurde geschrieben: \enquote{\textit{Rust is a solution that scales and is straightforward to deploy. It keeps resource usage low
        without the possibility of compromising memory safety. }}\autocite{rust-whitepaper-npm}
Dabei wird noch geschrieben, dass Rust langweilig sei, dass dies aber auch das beste Kompliment sei, was eine Programmiersprache erhalten k"onnte (\enquote{\textit{npm’s first Rust program hasn't
        caused any alerts in its year and a half
        in production.}}) \autocite{rust-whitepaper-npm}\\
Allerdings werden in dem Paper auch Nachteile durch den Einsatz von Rust bei npm genannt. Beispielsweise ist Rust eine junge Progarmmiersprache und dadurch hat diese noch keine Industrie-Standards f"ur Logging oder Alerting. Deshalb musste hierbei alles neu geschrieben werden, was mit viel Aufwand verbunden war.
% chapter Anwendungsbereiche (end)


% !TEX root =  master.tex
\chapter{Zusammenfassung}

\nocite{*}

Dieses Kapitel enthält die Zusammenfassung der Arbeit mit Fazit und Ausblick.

\section{Fazit}

In dieser Arbeit wurde zunächst ein Überblick zu den unterschiedlichen Kernkonzepten und Besonderheiten der Programmiersprache Rust wie Ownership, Borrowing, Lifetimes und unsicheren Code, der in sicheren APIs gekapselt ist, gegeben. Rust ist die erste von der Industrie unterstützte Sprache, die den jahrzehntelangen Kompromiss im Programmiersprachenentwurf zwischen Sicherheit und Kontrolle überwinden konnte. Sie ist sowohl eine typsichere Sprache als auch eine Systemprogrammiersprache. Hierbei unterscheidet sich Rust von den anderen Sprachen, die die gleiche Garantie für Typsicherheit bieten, darin, dass sie dies ohne Garbage Collector oder manuelle Speicherverwaltung tut. Rust ist nicht wirklich eine objektorientierte Sprache, obwohl sie einige objektorientierte Merkmale aufweist. Rust ist keine funktionale Sprache, obwohl sie dazu neigt, die Einflüsse auf das Ergebnis einer Berechnung deutlicher zu machen, wie es funktionale Sprachen tun. Rust ist für Entwickler und Projekte gedacht, bei denen nicht nur Leistung und low-level-Optimierungen wichtig sind, sondern bei denen auch eine sichere und stabile Ausführungsumgebung benötigt wird. Rust fügt eine Menge funktionaler High-Level-Programmiertechniken innerhalb der Sprache hinzu, so dass sie sich gleichzeitig wie eine Low-Level- und eine High-Level-Sprache anfühlt. Rust ähnelt C und \texttt{C++} bis zu einem gewissen Grad, aber viele Idiome aus diesen Sprachen treffen eher nicht zu, so dass der typische Rust-Code C oder \texttt{C++}-Code nicht sehr ähnlich ist. 

Rust ist für die Implementierung jener fundamentalen Systemschichten konzipiert, die eine starke Leistung und eine umfassende Kontrolle über die Ressourcen erfordern, aber dennoch das grundlegende Niveau der Vorhersagbarkeit garantieren, das die Typsicherheit bietet. Dadurch wird Entwicklern möglich, die Kosten ihres Codes vorherzusehen.
Rust hält den Ressourcenverbrauch niedrig, ohne hierbei, die Speichersicherheit zu beeinträchtigen. Rust bringt moderne Werkzeuge in den Bereich der Systemprogrammierung. Zwar gibt es immer noch Best Practices und Werkzeuge, die verbessert werden könnten, aber der Erfolg von Rust ist auf langfristigen Erfolg ausgerichtet.

\section{Ausblick}

...




\initializeAppendix

% !TEX root =  master.tex
\chapter{Beispiel-Anhang: Testanhang}
Anh\"ange\index{Anhang} werden am Ende Ihrer Arbeit vor dem Literaturverzeichnis und dem Index eingef\"ugt.

\section{Abschnitt im Anhang}


% !TEX root =  master.tex
\chapter{Beispiel-Anhang: Noch ein Testanhang}
nochmal: lipsum ...



\singlespacing


\initializeBibliography


\addcontentsline{toc}{chapter}{Index}
\printindex

\end{document}
